\documentclass[11pt,a4paper]{article}
\usepackage[T1]{fontenc}
\usepackage{lmodern,amsmath,amssymb,amsthm,mathtools,microtype}
\usepackage[margin=27mm]{geometry}
\usepackage[colorlinks=true,linkcolor=blue,citecolor=blue,urlcolor=blue]{hyperref}
\setlength{\emergencystretch}{2em}
\newtheorem{theorem}{Theorem}[section]
\newtheorem{lemma}[theorem]{Lemma}
\newtheorem{corollary}[theorem]{Corollary}
\newcommand{\Q}{\mathbb Q}\newcommand{\Z}{\mathbb Z}
\DeclareMathOperator{\Tr}{Tr}\DeclareMathOperator{\N}{N}
\title{Finiteness of automatic norm levels for conditional elliptic detectors}
\author{}\date{16 November 2026}
\begin{document}\maketitle
\begin{abstract}
For the conditional elliptic detectors, consider the complement level
$w\N(W)$ obtained from the identity involving the actual reciprocal
trace. At any positive odd contact valuation $\nu$, the two reference
tests force this level to have valuation at most $-\nu$. The assertion
holds for arbitrary guarded auxiliary points, at split and inert primes,
and does not require the norm of the detector to be a square.
Consequently an admissible automatic level confines the contact argument
to a finite set of rational numbers, by the finiteness theorem for
genus-two curves. Two such levels can occur at only finitely many
integer inputs for any fixed coefficient member, and at no input when
the fixed positive shift is sufficiently large. The result concerns
these automatic levels; it imposes no such conclusion on independently
chosen prescribed-trace norm incidences.
\end{abstract}

\section{The detector and its automatic complement level}

Put
\[
 K=\Q(\rho),\quad \rho^2=2,\quad D=11-6\rho,\quad
 E:\ y^2=x^3-Dx,
\]
and let
\[
 R_0=(1+\rho,2\rho),\qquad P=2R_0,
 \qquad x(P)=\frac{57}{8}-\frac72\rho.
\]
For a point in the affine domains below write
$f(Q)=1/x(Q)$ and $z(Q)=-x(Q)/y(Q)$.
All statements concern actual guarded tuples of the conditional source
\cite{Direct}. In particular, the required point coordinates,
differences and ratios are defined. Write $\sigma$ for conjugation and
set
\[
 L=\frac{z(U)}{z(H)},\qquad
 L'=\sigma\!\left(\frac{z(U-A)}{z(H-P)}\right).
\]
The six tests are
\begin{equation}\label{eq:tests}
 C_1=f(H),\quad C_2=f(U),\quad C_3=\sigma f(H-P),\quad
 C_4=\sigma f(U-A),\quad C_5=L-L',\quad C_6=L-X.
\end{equation}
Zero tests are allowed. Define
\[
 \Theta=\left\{0,-\frac12,-\frac14,-\frac25,-1,-\frac13\right\},
 \quad F(t)=\prod_{\theta\in\Theta}(t-\theta),
\]
\begin{equation}\label{eq:weights}
 \delta=t-1,\qquad
 w(t)=\frac{F(t)}{4\delta^8},\qquad
 \alpha(t)=\frac{\delta^4}{F(t)},\qquad
 W(t)=\prod_{i=1}^{6}(1+\alpha(t)C_i^2).
\end{equation}
The source guards impose $t\notin\Theta\cup\{1\}$ and
$\N(W)\Tr(W^{-1})\ne0$.
Its physical input is
\begin{equation}\label{eq:b}
                  b(t)=\frac{4w(t)}{\Tr(W(t)^{-1})^2}.
\end{equation}

Writing $W=W_0+W_1\rho$, define the \emph{automatic level}
\begin{equation}\label{eq:automatic}
                         \mathcal A(t)=w(t)\N(W(t)).
\end{equation}
The exact identity
\begin{equation}\label{eq:complement}
 2\mathcal A(t)(\mathcal A(t)-b(t))
 =\left(\frac{2w(t)\N(W(t))W_1}{W_0}\right)^2
\end{equation}
follows from $\N(W)=W_0^2-2W_1^2$; see~\cite{Benchmark}.
It supplies a rational complement but does not impose the required
arithmetic condition on its level.

For a finite set $S$ of rational primes, call a nonzero rational number
$a$ \emph{admissible outside $S$} if
\begin{equation}\label{eq:admissible}
                       v_p(a)\in2\Z_{\ge0}\qquad(p\notin S).
\end{equation}
The original source uses $S_0=\{p:p\le76176\}$.
The proofs below apply to every finite
\begin{equation}\label{eq:S}
                        S\supseteq\{2,3,5,7,41\}.
\end{equation}
They allow either sign of an admissible level and place no bound on
its size. In particular the conclusion is stronger than a restriction
to a fixed finite list of levels.

\section{A valuation bound for every auxiliary tuple}

\begin{lemma}\label{lem:contact}
Let $p\notin S$. If $v_p(w(t))$ is odd, then it is positive,
$v_p(\delta)=0$, and
\[
                v_p(w(t))=v_p(F(t))=:\nu>0,
                \qquad v_p(\alpha(t))=-\nu.
\]
\end{lemma}
\begin{proof}
The contacts are integral and pairwise distinct modulo $p$, and none
is congruent to $1$. All primes in the numerators and denominators
of their nonzero differences and their differences from $1$ belong
to $\{2,3,5,7\}$.
If $v_p(t)<0$, then
$v_p(F(t))=6v_p(t)$ and $v_p(\delta)=v_p(t)$, whence
$v_p(w(t))=-2v_p(t)$ is even. If $v_p(t)\ge0$ and
$v_p(\delta)>0$, every contact factor is a unit and
$v_p(w(t))=-8v_p(\delta)$ is even. Otherwise $\delta$ is a
unit and the contact valuation is nonnegative. This proves the
assertion, including the sign of the odd valuation.
\end{proof}

Normalize the valuations on $K$ at primes above an odd $p$ by
$v_{\mathfrak p}(p)=1$. They are integer-valued because $K/\Q$
is unramified away from $2$.

\begin{lemma}\label{lem:reference}
For every $p\notin S$ and every place $\mathfrak p\mid p$,
at least one of $f(H)$ and $f(H-P)$ has nonpositive valuation.
\end{lemma}
\begin{proof}
The displayed integral Weierstrass equation has good reduction:
its discriminant is $64D^3$ and $\N(D)=49$.
Furthermore
\[
                           \N(x(P))=\frac{41^2}{64},
\]
so $x(P)$ is a unit at both places above $p$, and the reduction of
$P$ is not the identity. A positive valuation of $f(Q)$ means that
$Q$ reduces to the identity. If both $H$ and $H-P$ reduced to the
identity, the reduction homomorphism would send their difference
$P$ to the identity, a contradiction.
\end{proof}

\begin{theorem}\label{thm:local}
For every actual guarded tuple and every $p\notin S$ with
$v_p(w(t))=\nu>0$ odd,
\begin{equation}\label{eq:bound}
                 v_p(\N(W(t)))\le-2\nu,
                 \qquad v_p(\mathcal A(t))\le-\nu<0.
\end{equation}
This holds whether $p$ splits or remains inert in $K$.
\end{theorem}
\begin{proof}
For a local value $c\in K$, allowing $c=0$, put
\[
              m_{\mathfrak p}(c)=\min\{0,2v_{\mathfrak p}(c)-\nu\},
              \qquad v_{\mathfrak p}(0)=+\infty.
\]
The two terms in $1+\alpha c^2$ have distinct valuations, since
$\nu$ is odd. Hence this factor is nonzero and has valuation
$m_{\mathfrak p}(c)\le0$.
Lemma~\ref{lem:reference} gives
\begin{equation}\label{eq:pairbound}
 m_{\mathfrak p}(f(H))+m_{\mathfrak p}(f(H-P))\le-\nu.
\end{equation}
Let $f_{\mathfrak p}$ denote the residue degree. All remaining
factors of $W$ have nonpositive valuation, and taking the norm
removes the conjugation in its third test. Therefore
\begin{align*}
 v_p(\N(W))
 &\le\sum_{\mathfrak p\mid p}f_{\mathfrak p}
       \bigl(m_{\mathfrak p}(f(H))+
             m_{\mathfrak p}(f(H-P))\bigr)\\
 &\le-\nu\sum_{\mathfrak p\mid p}f_{\mathfrak p}=-2\nu.
\end{align*}
In the split case this is a sum over two places of residue degree
one; in the inert case there is one place of residue degree two.
Adding $v_p(w)=\nu$ proves the second inequality.
\end{proof}

The theorem makes no assumption about the elliptic indices of
$A,H,U$, about a prime-selection procedure, or about the square
class of $\N(W)$. It uses the two reference tests already present
in~\eqref{eq:tests}. In particular it applies to both direct conditional
sources, with height-scale exponents four and six.

\section{Finiteness of admissible automatic levels}

\begin{theorem}\label{thm:finitearguments}
For every fixed finite $S$ satisfying~\eqref{eq:S}, there is a finite
set $\mathcal T_S\subset\Q\setminus(\Theta\cup\{1\})$ such that
every actual guarded tuple for which $\mathcal A(t)$ is admissible
outside $S$ has $t\in\mathcal T_S$.
The set is independent of all auxiliary points and elliptic indices.
\end{theorem}
\begin{proof}
If $v_p(w(t))$ were odd at any $p\notin S$,
Lemma~\ref{lem:contact} and Theorem~\ref{thm:local} would contradict
admissibility. Thus every outside-$S$ valuation of $w(t)$, and hence
of $F(t)$, is even.
The square class of $F(t)$ therefore belongs to the finite group
generated by $-1$ and the primes in $S$. Let $\mathcal D_S$ be the
set of signed squarefree representatives of this group. For some
$d\in\mathcal D_S$ there is a rational point on
\begin{equation}\label{eq:genus2}
                              y^2=dF(t).
\end{equation}
The polynomial $F$ has degree six and six distinct roots. Each
smooth projective model of~\eqref{eq:genus2} is therefore a
geometrically integral curve of genus two. Faltings' theorem
\cite[Satz~7]{Faltings} makes its set of rational points finite.
Taking the union of their finite rational $t$-coordinate sets,
and deleting the excluded arguments, gives $\mathcal T_S$.
No effective enumeration of this finite set is asserted.
\end{proof}

\begin{corollary}\label{cor:inputs}
Fix $d\in\{4,6\}$ and $B\ge3$. Suppose that the two arguments are
the actual coupled arguments of the conditional source,
\[
 t_0=\lambda,\qquad t_1=a\lambda,\qquad
 a=(X^2+B)^d+\bigl(\Tr x(A)\bigr)^d,
\]
and that both automatic levels are admissible outside the same fixed
finite $S$. Then only finitely many integer inputs $X$ are possible,
even when all auxiliary points and both admissible levels vary.
Moreover, when $B$ is sufficiently large in terms of $S$ and $d$,
there is no such guarded tuple at any rational input.
\end{corollary}
\begin{proof}
Both arguments belong to $\mathcal T_S$, and $t_0\ne0$ because
$0\in\Theta$ is guarded. Their positive ratio therefore belongs
to a fixed finite set. Let $M_S\ge1$ bound all those positive
ratios, taking $M_S=1$ if there are none. Since $d$ is even,
\[
                        (X^2+B)^d\le a\le M_S.
\]
For fixed $B$ this bounds $|X|$, hence gives only finitely many
integer inputs. If $B^d>M_S$, it is impossible for any rational $X$.
\end{proof}

\begin{corollary}\label{cor:scaling}
Fixed exceptional-unit rescalings of the physical weights and automatic
levels do not change either conclusion. More generally, a fixed finite
set of nonzero rational rescaling coefficients cannot yield automatic
admissible levels for infinitely many integer inputs.
\end{corollary}
\begin{proof}
An exceptional-unit rescaling leaves every outside-$S$ valuation
unchanged. For a finite set of arbitrary coefficients, let $S'$ be
$S$ together with their finite prime supports. Admissibility of a
rescaled level outside $S$ implies admissibility of the unscaled
level outside $S'$. Apply the preceding results to this one fixed
finite set $S'$. This enlargement is used only in the proof of
finiteness; it does not modify a source soundness predicate.
\end{proof}

These results do not exclude a finite prescribed-trace benchmark such
as that of~\cite{Benchmark}. They also do not concern levels chosen
independently of~\eqref{eq:automatic}, nor the redesigned general scalar
of~\cite{NormPreparation}. At a fixed native input $b\ne0$, the identity
\[
 a=\frac{b}{1-2u^2},\qquad q=\frac{2bu}{1-2u^2}
 \quad\Longrightarrow\quad q^2=2a(a-b)
\]
parametrizes rational level choices. The automatic level corresponds
to $u=W_1/W_0$. Choosing another value of $u$ supplies no automatic
admissibility or fixed-level theorem.

\section{Exact certificates}

The self-contained checker in
\nolinkurl{negative_results/computation/automatic_levels_v48}
verifies the reference constants, their local exceptional support, the
six simple contact roots, and the rational complement identities.
It recomputes actual detector values for two point configurations, at
one split prime and one inert prime outside $S_0$, with contact orders
one and three. Every example has the negative automatic-level
valuation in~\eqref{eq:bound}. It also checks that the two weights of
the prescribed-trace benchmark have no odd valuation outside $S_0$,
as required for consistency with that benchmark.
A finite valuation screen supplements the symbolic inequality
\eqref{eq:pairbound}; it is not used as a proof for all valuations.
The genus-two finiteness step is the cited theorem, not a computational
enumeration of rational points.

\begin{thebibliography}{9}
\bibitem{Direct}\emph{Thirteen Variables from Direct Elliptic Index Comparison},
\nolinkurl{papers/conditional_few_variables/direct_input13_20261110_v22.tex}.
\bibitem{Benchmark}\emph{Admissible complement levels and prescribed native traces},
\nolinkurl{papers/pole_parameters/admissible_levels_native_traces_20261116_v47.tex}.
\bibitem{NormPreparation}\emph{Norm preparation for general and conditional elliptic detectors},
\nolinkurl{papers/pole_parameters/norm_compatible_scalar_20261116_v47.tex}.
\bibitem{Faltings} G.~Faltings,
\emph{Endlichkeitss\"atze f\"ur abelsche Variet\"aten \"uber Zahlk\"orpern},
Invent. Math. \textbf{73} (1983), 349--366,
\url{https://doi.org/10.1007/BF01388432}; correction,
Invent. Math. \textbf{75} (1984), 381,
\url{https://doi.org/10.1007/BF01388572}.
\end{thebibliography}
\end{document}
